\subsubsection{Инициализация схем}

Схемы баз данных создаются заранее с помощью init-скриптов, а во время
эксперимента сервис создания нагрузки (loadgen) работает только с данными выбранной модели.

В PostgreSQL две модели размещены в одной базе, но в разных схемах:
\texttt{jsonb} и \texttt{normalized}. В MongoDB модели разделены по базам:
\texttt{coursework\_nested} и \texttt{coursework\_normalized}. Такое
разделение не смешивает данные разных моделей при последовательных запусках.

\subsubsection{Индексы}

Индексы подбирались под операции нагрузки. Для чтения за временной интервал
используются индексы по времени события и телеметрии:
\texttt{occurred\_at} и \texttt{recorded\_at}. Для выборок по месту
наблюдения добавлены составные индексы по территории, зоне и времени:
\texttt{area\_code}, \texttt{zone\_code}, \texttt{occurred\_at} или
\texttt{recorded\_at}. Для поиска важных событий используются индексы по
типу события, уровню значимости и времени.

В нормализованных моделях дополнительно индексируются поля связей:
\texttt{zone\_id}, \texttt{camera\_id}, \texttt{event\_id} и соответствующие
\texttt{\$ref}-ссылки в MongoDB. В PostgreSQL JSONB добавлены GIN-индексы по
\texttt{payload}, \texttt{settings} и \texttt{metrics}. В документных моделях
MongoDB индексируются вложенные поля, которые участвуют в запросах:
\texttt{payload.objects.object\_type}, \texttt{area.area\_code},
\texttt{zone.zone\_code}, \texttt{camera.serial\_number}.

\subsubsection{Подготовка данных}

Перед запуском нагрузки создается базовый мир видеонаблюдения: территории,
зоны, камеры, положения камер, настройки видеопотока, настройки аналитики,
области детекции и линии пересечения. Генерация детерминированная: одинаковые
\texttt{seed} и профиль размера дают одинаковые предметные данные для всех
моделей. События и телеметрия создаются уже во время нагрузочного цикла,
поскольку именно они образуют основной поток записей.

Для экспериментов предусмотрены профили \texttt{small}, \texttt{medium} и
\texttt{large}. В запуске из веб-интерфейса по умолчанию используется
\texttt{small}: 5 территорий, 5 зон на территорию и 5 камер на зону.
